%  LaTeX support: latex@mdpi.com 
%  For support, please attach all files needed for compiling as well as the log file, and specify your operating system, LaTeX version, and LaTeX editor.

%=================================================================
\documentclass[journal,article,submit,pdftex,moreauthors]{Definitions/mdpi}

%=================================================================
% MDPI internal commands - do not modify
\firstpage{1} 
\makeatletter 
\setcounter{page}{\@firstpage} 
\makeatother
\pubvolume{1}
\issuenum{1}
\articlenumber{0}
\pubyear{2026}
\copyrightyear{2026}
%\externaleditor{Firstname Lastname} % More than 1 editor, please add `` and '' before the last editor name
\datereceived{ } 
\daterevised{ } % Comment out if no revised date
\dateaccepted{ } 
\datepublished{ } 
%\datecorrected{} % For corrected papers: "Corrected: XXX" date in the original paper.
%\dateretracted{} % For retracted papers: "Retracted: XXX" date in the original paper.
%\doinum{} % Used for some special journals, like molbank
%\pdfoutput=1 % Uncommented for upload to arXiv.org
%\CorrStatement{yes}  % For updates
%\longauthorlist{yes} % For many authors that exceed the left citation part
%\IsAssociation{yes} % For association journals

%=================================================================
% Add packages and commands here. The following packages are loaded in our class file: fontenc, inputenc, calc, indentfirst, fancyhdr, graphicx, epstopdf, lastpage, ifthen, float, amsmath, amssymb, lineno, setspace, enumitem, mathpazo, booktabs, titlesec, etoolbox, tabto, xcolor, colortbl, soul, multirow, microtype, tikz, totcount, changepage, attrib, upgreek, array, tabularx, pbox, ragged2e, tocloft, marginnote, marginfix, enotez, amsthm, natbib, hyperref, cleveref, scrextend, url, geometry, newfloat, caption, draftwatermark, seqsplit
% cleveref: load \crefname definitions after \begin{document}

%=================================================================
% Please use the following mathematics environments: Theorem, Lemma, Corollary, Proposition, Characterization, Property, Problem, Example, ExamplesandDefinitions, Hypothesis, Remark, Definition, Notation, Assumption
%% For proofs, please use the proof environment (the amsthm package is loaded by the MDPI class).

%=================================================================
% Full title of the paper (Capitalized)
\Title{Evidence-Adaptive Chain-of-Thought Gating via Key-Fact Consistency for Budgeted Reliable Reasoning}

% Author Orchid ID: enter ID or remove command
\newcommand{\orcidauthorA}{0000-0000-0000-000X} % Add \orcidA{} behind the author's name
%\newcommand{\orcidauthorB}{0000-0000-0000-000X} % Add \orcidB{} behind the author's name

% Authors, for the paper (add full first names)
\Author{Firstname Lastname $^{1}$\orcidA{}, Firstname Lastname $^{2}$ and Firstname Lastname $^{2,}$*}

%\longauthorlist{yes}

% MDPI internal command: Authors, for metadata in PDF
\AuthorNames{Firstname Lastname, Firstname Lastname and Firstname Lastname}

% Affiliations / Addresses (Add [1] after \address if there is only one affiliation.)
\address{%
$^{1}$ \quad Affiliation 1; e-mail@e-mail.com\\
$^{2}$ \quad Affiliation 2; e-mail@e-mail.com}

% Contact information of the corresponding author
\corres{Correspondence: e-mail@e-mail.com; Tel.: (optional; include country code; if there are multiple corresponding authors, add author initials) +xx-xxxx-xxx-xxxx (F.L.)}

% Current address and/or shared authorship
%\firstnote{Current address: Affiliation.}
% Current address should not be the same as any items in the Affiliation section.

%\secondnote{These authors contributed equally to this work.}
% The commands \thirdnote{} till \eighthnote{} are available for further notes.

%\simplesumm{} % Simple summary

%\conference{} % An extended version of a conference paper

% Abstract (Do not insert blank lines, i.e. \\)
\abstract{Prompt-only adaptive chain-of-thought (CoT) methods often decide whether to spend extra reasoning tokens using a single self-reported confidence score. This signal is brittle: confidence is frequently miscalibrated, can be inflated by stylistic verbosity, and fails in the most socially harmful regime of confident but incorrect answers. We study how to allocate reasoning effort instance-by-instance using prompting alone while explicitly controlling overthinking and improving calibration-aware reliability rather than only average accuracy. We propose Evidence-Adaptive Chain-of-Thought (EA-CoT), a two-stage controller that augments confidence-only gating with factored metacognition. In stage 1, the model outputs a short draft answer without rationale, 2–4 atomic Key Facts that must hold for the answer to be correct, and confidence c. In stage 2, the same model tersely verifies each Key Fact against the question text and yields an evidence score e, the fraction marked supported. Full minimal CoT is triggered only if c is below a threshold or e is below a threshold. We specify a complete inference-only evaluation on the first 200 GSM8K test examples with a 50/150 tuning/evaluation split, comparing Direct, Fixed-CoT, confidence-adaptive CoT, and EA-CoT under matched token budgets. The provided artifact includes full code and methodology figures but not logged metric tables; accordingly, we contribute a rigorous protocol, metric definitions, and implementation analysis that determines whether EA-CoT reduces overconfident errors at comparable cost.}

% Keywords
\keyword{keyword 1; keyword 2; keyword 3 (List three to ten pertinent keywords specific to the article; yet reasonably common within the subject discipline.)}

% The fields PACS, MSC, and JEL may be left empty or commented out if not applicable
%\PACS{J0101}
%\MSC{}
%\JEL{}

%%%%%%%%%%%%%%%%%%%%%%%%%%%%%%%%%%%%%%%%%%
% Only for the journal Diversity
%\LSID{\url{http://}}

%%%%%%%%%%%%%%%%%%%%%%%%%%%%%%%%%%%%%%%%%%
% Only for the journal Applied Sciences
%\featuredapplication{Authors are encouraged to provide a concise description of the specific application or a potential application of the work. This section is not mandatory.}
%%%%%%%%%%%%%%%%%%%%%%%%%%%%%%%%%%%%%%%%%%

%%%%%%%%%%%%%%%%%%%%%%%%%%%%%%%%%%%%%%%%%%
% Only for the journal Data
%\dataset{DOI number or link to the deposited data set if the data set is published separately. If the data set shall be published as a supplement to this paper, this field will be filled by the journal editors. In this case, please submit the data set as a supplement.}
%\datasetlicense{License under which the data set is made available (CC0, CC-BY, CC-BY-SA, CC-BY-NC, etc.)}

%%%%%%%%%%%%%%%%%%%%%%%%%%%%%%%%%%%%%%%%%%
% Only for the journal BioTech, Fishes, Neuroimaging and Toxins
%\keycontribution{The breakthroughs or highlights of the manuscript. Authors can write one or two sentences to describe the most important part of the paper.}

%%%%%%%%%%%%%%%%%%%%%%%%%%%%%%%%%%%%%%%%%%
% Only for the journal Encyclopedia
%\encyclopediadef{For entry manuscripts only: please provide a brief overview of the entry title instead of an abstract.}

%%%%%%%%%%%%%%%%%%%%%%%%%%%%%%%%%%%%%%%%%%
% Different journals have different requirements. Please check the specific journal guidelines in the "Instructions for Authors" on the journal's official website.
%\addhighlights{yes}
%\renewcommand{\addhighlights}{%
%
%\noindent The goal is to increase the discoverability and readability of the article via search engines and other scholars. Highlights should not be a copy of the abstract, but a simple text allowing the reader to quickly and simplified find out what the article is about and what can be cited from it. Each of these parts should be devoted up to 2~bullet points.\vspace{3pt}\\
%\textbf{What are the main findings?}
% \begin{itemize}[labelsep=2.5mm,topsep=-3pt]
% \item First bullet.
% \item Second bullet.
% \end{itemize}\vspace{3pt}
%\textbf{What are the implications of the main findings?}
% \begin{itemize}[labelsep=2.5mm,topsep=-3pt]
% \item First bullet.
% \item Second bullet.
% \end{itemize}
%}

%%%%%%%%%%%%%%%%%%%%%%%%%%%%%%%%%%%%%%%%%%
\begin{document}

%%%%%%%%%%%%%%%%%%%%%%%%%%%%%%%%%%%%%%%%%%
\section{Introduction}

Large language models (LLMs) frequently improve problem-solving accuracy when prompted to produce intermediate reasoning, commonly referred to as chain-of-thought (CoT). However, always producing CoT has practical costs: it increases token usage and latency, creates user-facing verbosity, and can encourage overthinking on instances that are already easy. These tradeoffs motivate instance-adaptive prompting, where a system first runs a cheap pass and then decides whether to invest in deeper reasoning only when necessary.

A common adaptive design is a two-stage controller that asks the model for an answer and a self-reported confidence score, then triggers CoT when confidence is low. Instance-adaptive zero-shot CoT prompting formalizes this idea and shows that confidence-gated reasoning can retain much of CoT’s accuracy with fewer tokens on average in some settings [deutsch-2023-instance]. Yet confidence gating inherits a fundamental weakness: the control signal is a single scalar produced by the same generator whose outputs we are trying to regulate. Self-reported confidence can be miscalibrated, can correlate with superficial traits such as verbosity, and can shift substantially with prompt wording or model style. As a result, confidence-only controllers can skip reasoning precisely when it is most needed, producing confident but wrong answers that are particularly harmful for user trust and safety.

This paper targets that failure mode. We ask how to allocate reasoning effort instance-by-instance using only prompting while (i) reducing dependence on confidence calibration, (ii) explicitly limiting overthinking and verbosity, and (iii) improving calibration-aware reliability. Here, reliability is not only average accuracy, but also the propensity to avoid high-confidence errors and to deliver better accuracy under an explicit token budget.

We propose Evidence-Adaptive Chain-of-Thought (EA-CoT), a prompt-only controller that augments scalar confidence with a decomposed, verbalizable evidence signal. The motivating intuition mirrors a common human strategy: when an answer is genuinely easy, one can quickly state a few decisive facts that would have to be true for the answer to hold; when those facts are missing or inconsistent, it is a cue to slow down and reason more carefully. EA-CoT operationalizes this by asking the model to output a draft answer without rationale, a small set of answer-critical atomic Key Facts, and a confidence score. A second, intentionally terse call then checks whether each Key Fact is supported by the question text, yielding an evidence score e. Deep reasoning is triggered when either confidence is low or the evidence score is weak. This design introduces an orthogonal veto: high confidence alone is no longer sufficient to skip additional reasoning.

We instantiate and evaluate EA-CoT on GSM8K, a benchmark of grade-school math word problems with single numeric answers. The evaluation plan is inference-only and lightweight: we use the first 200 test items with a 50-item tuning split for selecting thresholds and a 150-item evaluation split for reporting metrics. We compare four methods: Direct (answer only), Fixed-CoT (always produce CoT), confidence-adaptive CoT (CA-CoT), and EA-CoT. The runner computes accuracy, average token usage as a cost proxy, and CoT trigger rate as a direct measure of how often deeper reasoning is invoked. It also computes a confidence-related overconfidence statistic labeled confident_wrong_rate.

The experimental artifact associated with this paper includes full runnable code, Hydra configurations, and two figures documenting methodology and system architecture. However, it does not include logged numerical outcomes (for example, no Weights and Biases summaries or exported per-run metric tables). Therefore, we focus on (a) a rigorous specification of the controller and evaluation protocol, and (b) an implementation analysis identifying details that determine whether EA-CoT can in fact improve calibration-aware reliability relative to confidence-only gating.

Contributions:
- We introduce EA-CoT, a prompt-only adaptive reasoning controller that uses both self-reported confidence c and an evidence-consistency score e computed from a Key-Fact decomposition.
- We specify a calibration-aware evaluation protocol for budgeted adaptive prompting on GSM8K, emphasizing matched-budget comparisons and overconfidence-oriented metrics.
- We provide a complete runnable codebase (inference runner, preprocessing, evaluation scripts, and Hydra configurations) for Direct, Fixed-CoT, CA-CoT, and EA-CoT on a fixed 200-example GSM8K slice.
- We identify implementation pitfalls that affect correctness and interpretability, including strict “question-only” evidence verification on math tasks, a substring-counting bug risk in evidence scoring, and a mismatch between the intended and implemented confident-wrong metrics.

Future work should (i) align the evidence-checking criterion with the structure of reasoning datasets (distinguishing facts stated in the prompt from facts that are legitimately derivable), (ii) correct evidence parsing and confidence-based reliability metrics, and (iii) report full accuracy–cost tradeoffs under explicit token budgets to quantify when evidence-based gating improves reliability beyond confidence-only controllers.

%%%%%%%%%%%%%%%%%%%%%%%%%%%%%%%%%%%%%%%%%%
\section{Related Work}

EA-CoT most directly builds on prompt-only instance-adaptive controllers that aim to reduce reasoning cost by invoking CoT selectively. Instance-adaptive zero-shot CoT prompting proposes a two-stage scheme that gates CoT using self-reported confidence, demonstrating improved accuracy–cost tradeoffs relative to always-on CoT in the studied settings [deutsch-2023-instance]. EA-CoT shares the same high-level controller template but changes the decision signal: instead of trusting only a scalar confidence c, it adds an explicit evidence decomposition (Key Facts) and an evidence-consistency score e derived from a constrained verification prompt. This difference is motivated by the observation that confidence is often miscalibrated or style-sensitive; an evidence gate can, in principle, catch cases where a model is confident but unable to state question-supported premises for its answer.

A second point of contrast is diagnostic richness. Confidence-only gating exposes only a scalar internal signal, limiting post-hoc analysis of failure modes. EA-CoT exposes both the Key Facts and the evidence score, enabling analyses of whether evidence scores separate correct from incorrect instances, which part of the gate triggered deeper reasoning, and whether low-evidence cases correspond to specific error types. While the current runner stores key_facts and evidence_score in per-example JSON outputs, it does not compute additional evidence-based statistics (such as AUC of e for predicting correctness). Thus, EA-CoT differs not only as a control policy but also in the observables it makes available for reliability analysis.

Several reference-candidate papers are not methodologically comparable because they do not address LLM prompting or inference-time control. “Reasoning Like Program Executors” studies elastic metrics and square-root velocity transforms for shape analysis, which is orthogonal to token-budgeted reasoning control in language models [navayazdani-2024-reasoning]. “CLeAR: Continual Learning on Algorithmic Reasoning for Human-like Intelligence” concerns mmWave networking for extended reality traffic, again unrelated to prompt-based metacognitive control [struye-2024-clear]. “Chain-of-Thought Reasoning Without Prompting” is a theoretical survey in algebraic combinatorics rather than an LLM reasoning method, and therefore cannot serve as a baseline for the present problem setting [rote-2023-chain]. We mention these works only to delimit scope and to emphasize that our comparisons are restricted to prompt-only inference-time controllers on question answering benchmarks.

%%%%%%%%%%%%%%%%%%%%%%%%%%%%%%%%%%%%%%%%%%
\section{Background}

We study inference-time control of a single black-box LLM accessed through prompting. Each dataset instance is a question x paired with a gold answer y. A controller may issue one or more prompts to the same model and returns a final prediction. No fine-tuning, external verifier model, tool use, or sampling-based self-consistency is employed; the constrained resource is tokens, used as a proxy for monetary cost, latency, and user-facing verbosity.

Problem setting and notation. Let M denote the language model as a text generator. A method defines a policy over a sequence of prompts p1(x), p2(x, r1), …, where rj is the text response at stage j. The method returns a final answer string â, which is converted into a task-specific prediction via an extraction function. For adaptive methods, the stage-1 response additionally provides internal signals: a scalar confidence c in [0, 1]; a set of Key Facts F = {f1, …, fk} with k between 2 and 4; and an evidence score e in [0, 1]. Each Key Fact is intended to be an atomic statement that must be true for the answer to be correct. The evidence score e is computed as the fraction of Key Facts labeled supported by a subsequent verification prompt.

Adaptive gating policies. A confidence-adaptive controller (CA-CoT) uses a two-stage template: (1) obtain a draft answer and confidence; (2) trigger a full CoT call if c < τc for a tuned threshold τc, otherwise return the draft answer. EA-CoT adds an evidence-gated condition: trigger CoT if (c < τc) or (e < τe), where τe is a second tuned threshold. The design intent is to reduce dependence on confidence calibration by adding a second gating signal grounded in a constrained evidence self-check.

Task instantiation: GSM8K and numeric answer extraction. The benchmark used in the artifact is the GSM8K test split. Evaluation uses the first 200 items with a 50/150 tuning/evaluation split. Correctness is exact numeric match to the gold final answer. Because prompts can elicit extra text and additional numbers (notably confidence values), the numeric extraction routine prioritizes certain patterns: it first matches a GSM8K-style “#### N”, then matches “Answer: N”, then matches “answer is N” (case-insensitive), and finally falls back to the last number in the response. This ordering is necessary because naive “last number” extraction can mistakenly return a confidence value rather than the answer.

Evaluation metrics and their implementation. The runner computes accuracy and average tokens per example, where token usage is summed across all calls used by a method. For the OpenAI client branch, the code uses total tokens (input plus output), making the metric a cost proxy rather than a pure verbosity measure; nevertheless it is consistent with the budgeted-control motivation. The runner also computes a CoT trigger rate. Finally, it computes a confidence-related statistic labeled confident_wrong_rate for adaptive methods. The intended definition from the experimental design is: among incorrect items, the fraction with c ≥ 0.8. The current implementation differs: it computes 1 minus the accuracy among items with confidence > 0.7. This implemented statistic measures the error rate within a high-confidence slice, not the fraction of errors that are high-confidence, and the difference affects interpretability.

Assumptions and constraints. All methods use the same model settings to isolate controller effects. Hydra configurations set temperature to 0.0 for near-deterministic decoding. Experiments are inference-only and sequential, with modest local compute requirements; runtime is dominated by remote model calls.

%%%%%%%%%%%%%%%%%%%%%%%%%%%%%%%%%%%%%%%%%%
\section{Method}

EA-CoT is a prompt-only controller that decides whether to invest in chain-of-thought reasoning on a per-instance basis. Its defining feature is to complement scalar confidence with a decomposed evidence signal intended to be harder to inflate by verbosity and more informative for reliability-centric gating.

Given a question x, EA-CoT runs up to three stages.

In stage 1, the controller requests a draft answer without rationale, a small set of Key Facts, and a confidence score. The prompt explicitly forbids showing reasoning and requires exactly three fields: Answer, KeyFacts, and Confidence. For GSM8K, the Answer field is constrained to be a final numeric answer. The KeyFacts field contains 2–4 short atomic facts that must be true for the answer to be correct. The controller parses the response into (i) a draft answer a1, (ii) a list of facts F, and (iii) a scalar confidence c. If parsing fails (for example, a1 is missing), this is treated as a trigger for deeper reasoning.

In stage 2, if any Key Facts were extracted, the controller issues a short verification prompt that asks the same model to judge each fact using only the question text. Each fact is labeled with a single token from {S, U, C}, denoting supported, unsupported, or unclear. The evidence score e is computed as the fraction of facts labeled supported. This stage is designed to be terse so that its overhead is small compared with a full CoT call.

The gating rule is: trigger the expensive reasoning stage if (c < τc) or (e < τe) or a1 is missing; otherwise, return a1. The thresholds τc and τe are tuned on a separate tuning split.

In stage 3, when triggered, the controller requests a solution using the fewest steps needed and then constrains the output format to exactly two final lines: a one-line sanity Check and the final Answer. The controller parses the final answer a3 and returns it.

Design rationale. EA-CoT targets the key failure mode of confidence-only gating: confident but wrong drafts. Under EA-CoT, a model cannot avoid deeper reasoning solely by reporting high confidence; it must also produce a coherent set of answer-critical facts that the verifier deems supported. Conversely, on easy instances where the model can produce consistent Key Facts and a satisfactory evidence score, the controller can skip CoT and control verbosity.

Implementation alignment and pitfalls in the provided codebase. The artifact contains both a generic reference implementation sketch and a runnable implementation in src/inference.py. The runnable version follows the same high-level structure but introduces two details that materially affect behavior. First, the verifier prompt asks whether facts are verifiable “from the question alone (not the solution)”; for GSM8K, many relevant premises are derived rather than stated, so this strict criterion may depress evidence_score and increase CoT triggering. Second, evidence_score is computed via substring counting of the token “SUPPORTED”; because “SUPPORTED” appears inside “UNSUPPORTED”, this can inflate the supported count, raising evidence_score and suppressing CoT triggering. Since gating depends directly on e, such parsing choices are part of the operational method definition and must be controlled for when interpreting results.

%%%%%%%%%%%%%%%%%%%%%%%%%%%%%%%%%%%%%%%%%%
\section{Experimental Setup}

We evaluate prompt-only inference-time controllers on GSM8K under a fixed inference-only protocol designed to isolate the effect of adaptive gating on the accuracy–cost tradeoff.

Dataset. The preprocessing module loads the GSM8K dataset (“main” configuration) and uses the test split. For a lightweight and scriptable experiment, it selects the first 200 test examples. Gold answers are parsed from GSM8K’s “####” delimiter and converted to float when possible.

Tuning and evaluation splits. There is no training. The first 50 examples form a tuning set for threshold selection; the remaining 150 form an evaluation set for reporting metrics.

Compared methods. Four methods are implemented.

Direct uses a single prompt that requests only the final numeric answer.

Fixed-CoT uses a single prompt that requests step-by-step reasoning followed by the final numeric answer.

CA-CoT uses two stages: stage 1 requests Answer and Confidence; if confidence is below a tuned threshold, stage 2 runs full CoT, otherwise the stage-1 answer is returned.

EA-CoT (proposed) uses three stages: stage 1 requests Answer, a fixed number of Key Facts (configured as 3 in Hydra), and Confidence; stage 2 verifies facts and computes evidence_score; stage 3 runs full CoT if either confidence or evidence are below tuned thresholds, otherwise the stage-1 answer is returned.

Model configuration and decoding. The Hydra configurations set provider to “openai” and the model name to “gpt-4o-mini”, with temperature 0.0. The maximum token limits are method-dependent: Direct uses max_tokens 50, while Fixed-CoT and the adaptive methods use max_tokens 500. In sanity_check mode, the main script increases max_tokens to 800 to reduce truncation risk for CoT outputs.

Threshold tuning. For computational speed, the implementation tunes thresholds using only the first 20 items of the tuning set. CA-CoT grid-searches confidence thresholds in {0.3, 0.4, 0.5, 0.6, 0.7, 0.8}. EA-CoT grid-searches confidence thresholds in {0.4, 0.5, 0.6, 0.7} and evidence thresholds in {0.3, 0.5, 0.7}. Each grid point is scored by accuracy on these 20 tuning examples, and the best thresholds are used for evaluation.

Metrics and logging. For each evaluation example, the runner stores the model response, extracted numeric answer, total tokens (summed across calls), whether CoT was triggered, and correctness. Aggregate metrics are accuracy, avg_tokens, cot_trigger_rate, and confident_wrong_rate (as implemented). When enabled, runs are logged to Weights and Biases; additionally, a JSON file is written locally containing configuration, tuned thresholds, aggregate metrics, and per-example outputs. A separate evaluation script can query W&B by run display name and export per-run and cross-run comparison plots.

%%%%%%%%%%%%%%%%%%%%%%%%%%%%%%%%%%%%%%%%%%
\section{Results}

The provided artifact includes runnable code and two saved figures, but it does not include numerical outcomes from executed runs (no Weights and Biases run summaries, no exported metrics tables, no aggregated comparison JSON, and no captured stdout/stderr). Therefore, the only reportable results with strict fidelity are: (i) the identity of the recorded figures, and (ii) the measurement definitions that the implementation is designed to compute.

![Methodology diagram for Evidence-Adaptive Chain-of-Thought (EA-CoT); this figure is conceptual rather than a performance plot, so “better” is not applicable.](images/ea_cot_methodology.pdf){#fig:ea-cot-methodology}

![System architecture diagram for the EA-CoT experimental codebase and evaluation flow; this figure is conceptual rather than a performance plot, so “better” is not applicable.](images/system_architecture.pdf){#fig:system-architecture}

Planned and implemented measurements. When executed, each run reports accuracy (higher is better), avg_tokens (lower is better as a cost proxy), cot_trigger_rate (lower generally indicates less overthinking, conditional on accuracy), and confident_wrong_rate (lower is better). For EA-CoT, per-example outputs also store evidence_score and key_facts, enabling further mechanism-level analyses even though these are not summarized in the current aggregate metric reporting.

Directly testable comparisons once logs exist. The evaluation protocol enables a matched-budget comparison between EA-CoT and CA-CoT by selecting thresholds that yield comparable avg_tokens and then comparing accuracy and overconfidence-oriented behavior. Additionally, per-example logs permit analysis of how often EA-CoT triggers CoT due to low evidence alone (e < τe with high c) versus low confidence, isolating the contribution of the evidence gate.

Implementation limitations that affect interpretability. The current code includes choices and potential bugs that can dominate outcomes. First, the evidence verifier is strict: it asks whether a Key Fact is directly verifiable from the question alone, which may systematically label derived but valid premises as unclear or unsupported in GSM8K, depressing evidence_score and increasing cot_trigger_rate. Second, evidence_score is computed by counting the substring “SUPPORTED”, which also appears inside “UNSUPPORTED”, potentially inflating evidence_score and suppressing CoT triggering. Third, the implemented confident_wrong_rate differs from the intended definition in the experimental design: it measures error rate among items with confidence > 0.7, not the fraction of errors with c ≥ 0.8. Fourth, threshold tuning uses only 20 of the 50 tuning examples, increasing variance in selected thresholds.

Within the constraints of this artifact, the principal results are thus the end-to-end method specification, the evaluation protocol, and the documentation of the system via @fig:ea-cot-methodology and @fig:system-architecture. Empirical claims about superiority over Direct, Fixed-CoT, or CA-CoT require the missing run summaries and exported metrics.

%%%%%%%%%%%%%%%%%%%%%%%%%%%%%%%%%%%%%%%%%%
\section{Discussion}



%%%%%%%%%%%%%%%%%%%%%%%%%%%%%%%%%%%%%%%%%%
\section{Conclusions}

We studied prompt-only instance-adaptive allocation of chain-of-thought reasoning under a token budget, motivated by the brittleness of scalar self-reported confidence as a gating signal. We proposed Evidence-Adaptive Chain-of-Thought (EA-CoT), a controller that augments confidence-based gating with a decomposed evidence representation: a small set of answer-critical Key Facts and a terse verification step that produces an evidence score e. The gating rule triggers minimal CoT only when confidence is low or evidence is weak, aiming to reduce overthinking while also reducing harmful confident-but-wrong behavior.

We specified a concrete inference-only evaluation on GSM8K (first 200 test items, 50/150 tuning/evaluation split) and documented a complete runnable codebase with Hydra configurations and an evaluation pipeline, along with methodology and system-architecture figures (@fig:ea-cot-methodology and @fig:system-architecture). Because numerical run summaries are not included in the provided artifact, the verified contributions in this paper are the controller definition, the budgeted evaluation protocol, and an implementation-level analysis of factors that determine whether evidence-based gating can deliver the intended calibration-aware reliability improvements.

Next steps are technically concrete: align the evidence verification criterion with the structure of reasoning benchmarks, fix evidence parsing so supported and unsupported are not conflated, implement the intended confident-wrong metric, tune thresholds on the full tuning set, and then report accuracy–cost tradeoffs under matched budgets. These steps would enable a definitive test of whether evidence-consistency gating improves reliability beyond confidence-only adaptive prompting.

%%%%%%%%%%%%%%%%%%%%%%%%%%%%%%%%%%%%%%%%%%
\vspace{6pt}

%%%%%%%%%%%%%%%%%%%%%%%%%%%%%%%%%%%%%%%%%%
\authorcontributions{For research articles with several authors, a short paragraph specifying their individual contributions must be provided. The following statements should be used ``Conceptualization, X.X. and Y.Y.; methodology, X.X.; software, X.X.; validation, X.X., Y.Y. and Z.Z.; formal analysis, X.X.; investigation, X.X.; resources, X.X.; data curation, X.X.; writing---original draft preparation, X.X.; writing---review and editing, X.X.; visualization, X.X.; supervision, X.X.; project administration, X.X.; funding acquisition, Y.Y. All authors have read and agreed to the published version of the manuscript.'', please turn to the  \href{http://img.mdpi.org/data/contributor-role-instruction.pdf}{CRediT taxonomy} for the term explanation. Authorship must be limited to those who have contributed substantially to the work~reported.}

\funding{This research received no external funding.}

\dataavailability{All resources used in this study are openly available at }

\acknowledgments{In this study, we automatically carried out a series of research processes—from hypothesis formulation to paper writing—using generative AI.}

\conflictsofinterest{The authors declare no conflicts of interest.}

%%%%%%%%%%%%%%%%%%%%%%%%%%%%%%%%%%%%%%%%%%
%\isPreprints{}{% This command is only used for ``preprints''.
\begin{adjustwidth}{-\extralength}{0cm}
%} % If the paper is ``preprints'', please uncomment this parenthesis.
%\printendnotes[custom] % Un-comment to print a list of endnotes

\bibliographystyle{plainnat}
\bibliography{references}

\PublishersNote{}
%\isPreprints{}{% This command is only used for ``preprints''.
\end{adjustwidth}
%} % If the paper is ``preprints'', please uncomment this parenthesis.
\end{document}
